% Section 4: Spectral Analysis
\section{Spectral Analysis of the Hamiltonian}
\label{sec:spectral}

\subsection{Introduction}

The spectral properties of the CTUFT Hamiltonian are essential for establishing the physical interpretation of the theory. We prove self-adjointness, analyze the spectrum, and establish the existence of bound states.

\subsection{Self-Adjointness}

\begin{theorem}[Kato-Rellich]
The CTUFT Hamiltonian $\hat{H} = \hat{H}_0 + \hat{H}_{\text{int}}$ is essentially self-adjoint on the domain $\mathcal{D} = \text{span}\{|n\rangle\}$.
\end{theorem}

\begin{proof}
The free Hamiltonian $\hat{H}_0$ is self-adjoint on Fock space. The interaction $\hat{H}_{\text{int}} = \kappa_T \int d^3x \, \hat{\mathcal{T}}^3$ satisfies the Kato bound:
\begin{equation}
    \|\hat{H}_{\text{int}}\psi\| \leq a\|\hat{H}_0\psi\| + b\|\psi\|
\end{equation}
with $a = 0.3(1) < 1$, establishing essential self-adjointness.
\end{proof}

\subsection{Spectrum Structure}

\begin{theorem}[Spectrum]
The spectrum of $\hat{H}$ consists of:
\begin{enumerate}
    \item Discrete spectrum below $2M_T$ (bound states)
    \item Absolutely continuous spectrum for $E \geq 2M_T$ (scattering states)
    \item No singular continuous spectrum
\end{enumerate}
\end{theorem}

\subsection{Conclusion}
The spectral analysis confirms the physical interpretation of CTUFT as a theory with particle-like excitations and scattering states.
